\documentclass[11pt]{article}
\usepackage[a4paper,margin=27mm]{geometry}
\usepackage{amsmath,amssymb,amsthm,booktabs,array,longtable}
\usepackage[T1]{fontenc}
\usepackage{lmodern}
\usepackage{xurl}
\usepackage[colorlinks=true,linkcolor=blue,citecolor=blue,urlcolor=blue]{hyperref}
\newtheorem{theorem}{Theorem}
\newtheorem{lemma}{Lemma}
\newtheorem{corollary}{Corollary}
\newcommand{\ZZ}{\mathbb Z}
\newcommand{\QQ}{\mathbb Q}
\newcommand{\FF}{\mathbb F}
\newcommand{\coeff}[2]{[#1]#2}
\DeclareMathOperator{\End}{End}
\title{An elementary proof of Su's theorem on Bala's\newline inverse-function period conjectures,\newline with odd-exponent extensions}
\author{Haobo Ma\\\small ChronoAI Pte Ltd, Singapore
\and Wenlin Zhang\\\small National University of Singapore\\\small ChronoAI Pte Ltd, Singapore}
\date{14 November 2026}
\begin{document}
\maketitle
\begin{abstract}
We give an alternative proof of Lien-Hung Su's theorem (31 October 2026, Theorem 3.2 and Corollary 3.3 \cite{su}) resolving Peter Bala's totient-square period conjecture in OEIS A281440 (COMMENT, 19 February 2026; revision 17, 6 November 2026) \cite{oeis440}.
For the integral coefficients $a(n)=(2n)![x^{2n}]C(x)$ defined by
$C(x)=\cosh(\int_0^x C(t)^7\,dt)$, every modulus $m\ge3$ admits the eventual period $\varphi(m)^2/2$; moduli $1,2$ admit period one.
The proof replaces idempotent lifting by an elementary binomial identity for endomorphisms and includes the vanishing branch at the prime five.
The same argument covers the four companion conjectures in A281180, A281181, A281183 and A281184, and extends the half-square bound to every integer exponent $a\ge2$ in the associated quadratic derivation.
We prove the resulting odd-exponent corollaries, source derivative and algebraic identities, and finite-dimensional certificates for the exact period $72$ modulo $13$ and the two displayed A281428 examples.
\end{abstract}

\section{The source sequence and the theorem}
Bala asks whether the sequence in A281440 has an eventual period dividing
$\varphi(m)^2$ for every positive integer $m$ \cite{oeis440}.
Here $\varphi$ is Euler's totient function, and an eventual period $P$ means
$a(n+P)\equiv a(n)\pmod m$ for all sufficiently large $n$.
Su proved the stronger bound $\varphi(m)^2/2$ for $m\ge3$ within a uniform theorem
for inverse-function sequences \cite[Theorem 3.2]{su}.
We prove that theorem below using an elementary lifting lemma, then match its
parameters to the source equations. The congruence results and the seven
source-specific least periods in Section~\ref{sec:finite} are already in Su's
Theorem 3.2, Corollary 3.3 and Proposition 3.6.
Corollary~\ref{cor:odd} extends the argument to odd exponents; the separate
totient (rather than half-square) bound for A281428 remains Su's Theorem 4.1.

For an integer $a\ge2$ define $S_a,C_a\in\QQ[[x]]$ by
\begin{equation}\label{eq:flow}
 S_a'=C_a^a,\qquad C_a'=S_aC_a^{a-1},\qquad S_a(0)=0,\quad C_a(0)=1.
\end{equation}
Coefficient recursion gives a unique solution. Differentiating $C_a^2-S_a^2$
shows that it equals one. If $H_a(0)=0$ and $H_a'=C_a^{a-1}$, then
$(C_a+S_a)'=H_a'(C_a+S_a)$ and
$(C_a-S_a)'=-H_a'(C_a-S_a)$. Uniqueness for these linear formal differential
equations gives $C_a=\cosh H_a$ and $S_a=\sinh H_a$.
Conversely, $C=\cosh H$, $H'=C^{a-1}$, $H(0)=0$, gives \eqref{eq:flow} with
$S=\sinh H$. Thus A281440 is exactly the $a=8$ flow, with the even indexing
used in its source, not an ordinary-coefficient normalization.

Use the integral algebra and derivation
\begin{equation}\label{eq:derivation}
 R=\ZZ[s,c]/(c^2-s^2-1),\qquad
 D(s)=c^a,\quad D(c)=sc^{a-1},\quad E(f)=f(0,1).
\end{equation}
The relation is preserved since $D(c^2-s^2-1)=0$. The exponential flow
$\sum_{j\ge0}E(D^jf)x^j/j!$ is multiplicative by the iterated Leibniz rule.
Applied to $s,c$, it therefore gives precisely \eqref{eq:flow}.
For $k,n\ge0$ put
\begin{equation}\label{eq:columns}
 b_{a,k}(n)=E(D^{2n}c^k)=(2n)![x^{2n}]C_a(x)^k\in\ZZ.
\end{equation}
This proves integrality as well as the coefficient identification.
For direct integer evaluation one may also use
\begin{equation}\label{eq:recurrence}
 b_{a,k}(0)=1,\quad
 b_{a,k}(n)=k\bigl((k+a-1)b_{a,k+2a-2}(n-1)
              -(k+a-2)b_{a,k+2a-4}(n-1)\bigr).
\end{equation}
Indeed $D^2c^k=k\bigl((k+a-1)c^{k+2a-2}-(k+a-2)c^{k+2a-4}\bigr)$;
apply $E D^{2n-2}$. The $a=4$ recurrence is credited to Seiichi Manyama
in the source \cite{oeis181}; Su states the general form
\cite[Lemma 3.1]{su}.

\begin{theorem}[Su's even-exponent period theorem]\label{thm:period}
Let $a=2d\ge2$, $k\ge0$, and $m\ge3$. For each odd prime $p$, put
\[
 \alpha_{a,p}=(-1)^h\binom{d+h-1}{h}\pmod p,\qquad h=(p-1)/2.
\]
For $p^r\Vert m$, define
\[
 N_a(p^r)=
 \begin{cases}
 r,&p=2\text{ or }(p\text{ odd and }\alpha_{a,p}\ne0),\\
 \left\lceil r\bigl(1+(p-1)p^{r-1}\bigr)/2\right\rceil,
      &p\text{ odd and }\alpha_{a,p}=0,
 \end{cases}
\]
Put $N_a(m)=\max_{p^r\Vert m}N_a(p^r)$. Then for every $n\ge N_a(m)$,
\begin{equation}\label{eq:uniform}
 b_{a,k}\bigl(n+\varphi(m)^2/2\bigr)\equiv b_{a,k}(n)\pmod m.
\end{equation}
For $m=1,2$ an eventual period of one is valid.
\end{theorem}

\section{An elementary prime-power proof}
\begin{lemma}\label{lem:lift}
Let $M$ be a finite free module over $\ZZ/p^r\ZZ$, let $A\in\End(M)$, and
suppose $A^{T+1}\equiv A\pmod p$, with $T\ge1$. Then
\[
 (A^{Tp^{r-1}}-I)A^r=0.
\]
\end{lemma}
\begin{proof}
Set $B=A^T-I$ and $L=p^{r-1}$. Since $BA$ is divisible by $p$, the commuting
product $B^jA^r$ is divisible by $p^j$ when $1\le j\le r$, and by $p^r$
when $j>r$. Expand
\[
 (A^{TL}-I)A^r=\sum_{j=1}^L\binom Lj B^jA^r.
\]
For $1\le j\le\min(r,L)$, the identity
$j\binom Lj=L\binom{L-1}{j-1}$ gives
$v_p(\binom Lj)\ge r-1-v_p(j)$. Since $j-v_p(j)\ge1$, each summand is
divisible by $p^r$. The remaining summands already contain $p^r$.
\end{proof}

Fix a prime power $q=p^r$. The algebra
\[
 M_q=(\ZZ/q\ZZ)[s,c]/(s^q,c^2-s^2-1)
\]
is free of rank $2q$, with basis $s^j,cs^j$ for $0\le j<q$.
The derivation descends because $D(s^q)=qs^{q-1}c^a=0$. Evaluation at
$(0,1)$ remains well defined, so endomorphism identities on this algebra give
the desired coefficient congruences.

Suppose first that $p$ is odd. Modulo $p$ the polynomial $1+s^2$ has a unique
square root $u=1+O(s^2)$ in $\FF_p[s]/(s^q)$, obtained recursively using the
invertibility of two. The factors $c-u,c+u$ are comaximal because $2u$ is a
unit. On either resulting branch, $c^a=(1+s^2)^d$ since $a=2d$, and
\[
 D=f(s)\frac{d}{ds},\qquad f(s)=(1+s^2)^d.
\]
The $p$th iterate of a derivation is a derivation: the intermediate binomial
coefficients in its product rule vanish in characteristic $p$.
Every derivation on this truncated polynomial algebra is determined by its
value at $s$. Since $f$ is a unit, $D^p=\lambda D$ for some $\lambda$.
Commutation with $D$ gives $D\lambda=0$, whence
$\lambda\in\FF_p[s^p]/(s^q)$.

We compute its constant term explicitly. In $\FF_p[s]/(s^p)$ write
$f^{-1}=\sum_{j=0}^{p-1}\beta_js^j$ and put
$y=\sum_{j=0}^{p-2}\beta_js^{j+1}/(j+1)$.
Here $\beta_0=1$, so $y=s+O(s^2)$ is a coordinate. If
$\alpha=\beta_{p-1}$, then
\[
 Dy=1-\alpha s^{p-1}f=1-\alpha y^{p-1}.
\]
The scalar in $D^p=\lambda D$ is constant in this smaller quotient.
Evaluating $D^py=-\alpha D^{p-1}y^{p-1}$ at zero gives
$-\alpha(p-1)!=\alpha$ by Wilson's theorem. Thus
\[
 \lambda(0)=\alpha=[s^{p-1}](1+s^2)^{-d}
     =(-1)^h\binom{d+h-1}{h}=\alpha_{a,p}.
\]
On the larger quotient, $D\lambda=0$ implies inductively
$D^{1+(p-1)t}=\lambda^tD$. Frobenius and $s^{p^r}=0$ give
\begin{equation}\label{eq:restricted}
 D^J=\alpha D\pmod p,\qquad J=1+(p-1)p^{r-1}.
\end{equation}
If $\alpha=0$, the endomorphism $D^J$ is divisible by $p$, so
$D^{rJ}=0\pmod{p^r}$. Hence every $b_{a,k}(n)$ vanishes modulo $p^r$
for $2n\ge rJ$, and every positive shift is then a period.

If $\alpha\ne0$, take $t=(p-1)p^{r-1}$ in the preceding identity.
It gives $D^{T+1}=D\pmod p$ for $T=(p-1)^2p^{r-1}$.
Lemma~\ref{lem:lift}, with $A=D$, proves
\[
 (D^{\varphi(p^r)^2}-I)D^r=0\pmod{p^r}.
\]
Since $\varphi(p^r)^2$ is even, evaluation on $D^{2n}c^k$ gives the
column-index shift $\varphi(p^r)^2/2$ whenever $n\ge r$.

At $p=2$, let $A=D^2$. Modulo two it is itself a derivation, and direct
differentiation of \eqref{eq:derivation} gives
\[
 A(s)=0,\qquad A(c)=c(1+s^2)^{a-2}.
\]
Thus on $u(s)+cv(s)$, $A$ kills the first summand and multiplies the second
by $w=(1+s^2)^{a-2}$. In the quotient $s^{2^r}=0$, for $r\ge2$,
\[
 w^{2^{r-2}}=(1+s^{2^{r-1}})^{a-2}=1\pmod2,
\]
because $a-2$ is even and the square of $s^{2^{r-1}}$ vanishes.
For $r=1$ take $T=1$; for $r\ge2$ take $T=2^{r-2}$.
Then $A^{T+1}=A\pmod2$, and Lemma~\ref{lem:lift} yields a period
$T2^{r-1}$ in the column index once $n\ge r$.
This is $\varphi(2^r)^2/2$ for $r\ge2$, and one for $r=1$.

Finally, $\varphi(p^r)$ divides $\varphi(m)$ for every $p^r\Vert m$.
For odd $p$, and for $p=2,r\ge2$, both totients are even, so the local
half-square period divides $\varphi(m)^2/2$. The remaining local period
is one. Vanishing branches admit every shift. The Chinese remainder
theorem proves \eqref{eq:uniform}, including the stated threshold.
Modulus one is immediate, and the $r=1$ two-primary argument proves
the assertion for modulus two. This completes the proof of
Theorem~\ref{thm:period}.

For A281440 the scalar is especially transparent:
\[
 \alpha_{8,3}=2,\qquad \alpha_{8,5}=0,\qquad
 \alpha_{8,p}=(-1)^{(p-1)/2}\frac5{16}\quad(p\ge7).
\]
Thus five is its only odd vanishing prime. In particular,
$a(n)\equiv0\pmod{5^r}$ once
$n\ge\lceil r(1+4\cdot5^{r-1})/2\rceil$.
For $a=4$, $\alpha_{4,p}=(-1)^{(p-1)/2}/2\ne0$ at every odd prime.

\begin{corollary}[All integer exponents]\label{cor:odd}
Theorem~\ref{thm:period} holds for every integer $a\ge2$, using
\[
 \alpha_{a,p}=(-1)^h\prod_{j=0}^{h-1}\frac{a+2j}{2(j+1)}\in\FF_p,
 \qquad h=(p-1)/2.
\]
For odd $a\ge3$, this scalar vanishes at every prime $p\ge a$.
For the odd source flows $a=5,7,9$, the only nonvanishing odd primes
are respectively $\{3\}$, $\{3,5\}$ and $\{5,7\}$.
Consequently all their columns eventually vanish modulo every positive
integer coprime to $6$, $30$ and $70$, respectively.
\end{corollary}
\begin{proof}
For odd $p$, use the same square root $u^2=1+s^2$ in $\FF_p[s]/(s^{p^r})$.
The two branches are now $D=\pm u^a\partial_s$.
The coordinate calculation above applies to the unit $f=u^a$; the coefficient
$[s^{p-1}]f^{-1}$ is the displayed product, whose denominators are units.
Replacing $D$ by $-D$ leaves the scalar in $D^p=\lambda D$ unchanged because
$p$ is odd. Thus $D\lambda=0$, the Frobenius argument and both odd-prime
lifting cases are identical to those above.

At two put $A=D^2$. For odd $a$, direct differentiation gives
$A(s)=sw$ and $A(c)=cw$ modulo two, with $w=(1+s^2)^{a-1}$ and $A(w)=0$.
Therefore
$A(s^jc^e)=(j+e)w s^jc^e$ for $e=0,1$.
For $r=1$ take $T=1$; for $r\ge2$ take $T=2^{r-2}$.
Since $a-1$ is even,
$w^T=(1+s^{2^{r-1}})^{a-1}=1$ in the latter case; for $r=1$ also $w=1$.
On both parity classes this gives $A^{T+1}=A$ modulo two.
Lemma~\ref{lem:lift} supplies exactly the same two-primary shift and threshold.
The Chinese remainder argument is unchanged.

If $a$ is odd and $p\ge a$, the numerator product contains $p$, at
$j=(p-a)/2$, whereas its denominator is prime to $p$.
Checking the odd primes below $5,7,9$ in the product gives the stated sets.
The vanishing branch and the Chinese remainder theorem prove the final assertion.
\end{proof}

\section{Companion statements and source coverage}
The identities below follow from \eqref{eq:flow} and the source definitions
\cite{oeis180,oeis181,oeis183,oeis184,oeis440,oeiscompanions}:
\begin{center}
\begin{tabular}{lll}
\toprule
Entry & Column & Index range\\
\midrule
A281180 & $b_{4,4}(n-1)$ & $n\ge1$\\
A281181 & $b_{4,1}(n)$ & $n\ge0$\\
A281183 & $b_{4,2}(n)$ & $n\ge0$\\
A281184 & $b_{4,3}(n)$ & $n\ge0$\\
A281440 & $b_{8,1}(n)$ & $n\ge0$\\
A281439 & $b_{8,8}(n-1)$ & $n\ge1$\\
A281432 & $b_{6,1}(n)$ & $n\ge0$\\
A281427 & $b_{5,5}(n-1)$ & $n\ge1$\\
A281428 & $b_{5,1}(n)$ & $n\ge0$\\
A281436 & $b_{7,1}(n)$ & $n\ge0$\\
A281443 & $b_{9,9}(n-1)$ & $n\ge1$\\
A281444 & $b_{9,1}(n)$ & $n\ge0$\\
\bottomrule
\end{tabular}
\end{center}
For the odd series, differentiating $S_a$ gives
$(2n+1)![x^{2n+1}]S_a=b_{a,a}(n)$.
A281180, A281427, A281439 and A281443 list these coefficients from OEIS
index one. Their first value is $S_a'(0)=1$, so their column index is $n-1$.
\begin{corollary}[Source periods]\label{cor:sources}
The five totient-square conjectures in the first five rows are proved by
Theorem~\ref{thm:period}, as in Su's Corollary 3.3.
Every row has eventual period $\varphi(m)^2/2$ for $m\ge3$, and period one
for $m=1,2$. For A281180, A281427, A281439 and A281443 add one to the column
threshold; all other rows use $N_a(m)$ directly.
The odd-exponent half-square statements are consequences of
Corollary~\ref{cor:odd}, not conjectures asserted by those entries.
\end{corollary}
\begin{proof}
The column identifications follow from the source definitions and $S_a'=C_a^a$.
Apply Theorem~\ref{thm:period} and Corollary~\ref{cor:odd}.
\end{proof}

\begin{corollary}[Source derivative identities]\label{cor:derivatives}
For any $a\ge2$, let $\theta(0)=0$, $\theta'=\sec^{a-2}\theta$.
Then
\begin{equation}\label{eq:trig}
 (2n)![x^{2n}]C_a(x)
   =\left.\left(\sec^{a-2}\theta\,\frac d{d\theta}\right)^{2n}
       \sec\theta\right|_{\theta=0}.
\end{equation}
The formula marked with a question mark in A281436
(Bala, FORMULA, 23 December 2025), with $a=7$, as well as the displayed
formulas in A281432, A281440 and A281444 are therefore valid.
The tentative A281183 formula is valid as well.
\end{corollary}
\begin{proof}
The substitutions $s=\tan\theta$, $c=\sec\theta$ satisfy \eqref{eq:flow}.
Uniqueness and the chain rule prove \eqref{eq:trig}.
For A281183 put $u(z)=1+z^2$, $D_0=u^2\partial_z$ and
$L f=u\partial_z(uf)$. Then $L=u^{-1}D_0u$, so
$L^j1=u^{-1}D_0^ju$. Evaluation at zero and the $a=4$ flow give
$L^{2n}1(0)=b_{4,2}(n)$. This is the tentative derivative formula in
its FORMULA block, already proved by Su \cite[Proposition 3.5]{su}.
\end{proof}

These derivative representations belong to the inverse-function methods of
Dominici \cite[Theorem 4.1]{dominici} and Bala \cite[Corollary 1]{graves}.
Their other nested-derivative instances follow from a single conjugation:
for a formal unit $F(v)$ with $F(0)=1$, put $V=F\partial_v$ and
$L=\partial_v\circ F$. Then
\begin{equation}\label{eq:leftnested}
 L=F^{-1}VF,\qquad (L^jg)(0)=(V^j(Fg))(0).
\end{equation}
If $v'=F(v)$, $v(0)=0$, the right side is $j![x^j](Fg)(v(x))$.
For A281428 take $F=\sec^3v$, $g=\cos^2v$; this proves its recorded
left-nested formula since $Fg=\sec v$ and the flow gives $C_5$.
For A281181 take $F=(1-v^2)^{-1/2}$, $g=(1-v^2)^{-k/2}$;
the result is $(2n)![x^{2n}]C_4^{k+1}$, specifying the exponential
coefficient convention for its parameterized formula.

\begin{corollary}[Algebraic companion identities]\label{cor:algebraic}
For $a=2d+1\ge3$ set $T=S_a/C_a$. Then
\begin{equation}\label{eq:algebraic}
 x=\sum_{j=0}^{d-1}\frac{(-1)^j\binom{d-1}{j}}{2j+1}T^{2j+1},
 \qquad C_a=(1-T^2)^{-1/2},\quad S_a=TC_a.
\end{equation}
In particular, the algebraic equations for $S_5,C_5,T$ in A281427
(Bala, FORMULA, 14 December 2025) and $C_7$ in A281436
(FORMULA, 23 December 2025) hold. In A281181 the displayed term
\texttt{13*x\^{}5/6!} in $T=S_4/C_4$ (FORMULA, 14 December 2025)
must instead be $13x^5/5!$.
\end{corollary}
\begin{proof}
The quotient rule gives $T'=C_a^{a-3}$ and $1-T^2=C_a^{-2}$.
Thus $dx/dT=(1-T^2)^{d-1}$, whose integral is \eqref{eq:algebraic}.
For $a=5$ this is $3x=3T-T^3$. Substitution gives
\[
 (4-9x^2)C_5^6-3C_5^2-1=0,
\]
and, since $C_5^2=1+S_5^2$, also
\[
 (4-9x^2)S_5^6+3(4-9x^2)S_5^4+9(1-3x^2)S_5^2-9x^2=0.
\]
For $a=7$, $15x=15T-10T^3+3T^5$; inserting
$C_7^2=(1-T^2)^{-1}$ gives
$(225x^2-64)C_7^{10}+40C_7^4+15C_7^2+9=0$.
All branches have the initial values of \eqref{eq:flow}.
Finally $T'=C_4$ in the $a=4$ case, and
$C_4=1+x^2/2!+13x^4/4!+\cdots$ by \eqref{eq:recurrence};
integration proves the corrected $x^5$ coefficient.
\end{proof}

{\small
\begin{longtable}{p{.25\linewidth}p{.66\linewidth}}
\toprule
Source item & Coverage\\
\midrule
Five conjectures above & Proved: alternative proof of Su's Theorem 3.2 and Corollary 3.3, with the exact source columns and offsets.\\
A281439, A281432 & Not posed; period bound proved as companion columns.\\
A281440 mod-$13$ example & Proved: least period and scalar-column structure in Section~\ref{sec:finite}; prior proof in Su's Proposition 3.6.\\
A281181/183/184 residue examples & Proved: exact periods and mod-$12$ identity below; prior proofs in Su's Proposition 3.6.\\
A281183 derivative formula & Proved by Corollary~\ref{cor:derivatives}; prior proof in Su's Proposition 3.5.\\
A281436 derivative question & Answered: yes, by Corollary~\ref{cor:derivatives}; the question is the final FORMULA of 23 December 2025.\\
A281428 conjectures & Prior proof cited: Su's Theorem 4.1 gives eventual zero away from $2,3$, and eventual period dividing $\varphi(m)$; not reproved here.\\
A281428 two examples & Proved: zero modulo $385$ from index $6$, and least period $18$ modulo $27$ from index $7$, by Section~\ref{sec:finite}.\\
A281427, A281436, A281443, A281444 & Not posed; odd-exponent half-square bounds proved in Corollaries~\ref{cor:odd}--\ref{cor:sources}. Vanishing is proved away from $6,30,70$ for $a=5,7,9$, respectively.\\
A281427 and A281436 algebraic formulas & Proved as recorded, by Corollary~\ref{cor:algebraic}; these formulas predate the period conjectures.\\
A281181 displayed $T$ series & Refuted-and-corrected: the term printed as \texttt{13*x\^{}5/6!} is $13x^5/5!$ (FORMULA, 14 December 2025).\\
Nested-derivative formulas & Proved via \eqref{eq:trig} and \eqref{eq:leftnested}, with the exponential coefficient normalization. Earlier methods: Dominici and Bala.\\
General rectangular-layout remarks & The prime scalar relation is proved below; Corollary~\ref{cor:rows} gives Su's higher-order composite row relation for $a=4$. A universal rank-one interpretation at every composite modulus is not asserted in the source and is not adopted here.\\
\bottomrule
\end{longtable}}

\section{Finite-dimensional certificates}\label{sec:finite}
Here is a finite proof of the source's mod-$13$ example, rather than an
extrapolation from a long prefix. Write $z=s^2$ and $c^ef(z)$, where
$e\in\{0,1\}$. On this parity component $D^2$ induces
\begin{equation}\label{eq:matrix}
 g=ef+2(1+z)f',\qquad
 B_{a,e}f=(1+z)^{a-2}\bigl((1+(e+a-1)z)g+2z(1+z)g'\bigr).
\end{equation}
For an integer $m\ge2$, put $h=\lceil m/2\rceil$ and reduce modulo
$(m,z^h)$. The only degree-lowering coefficient on $z^h$ is
$2h(2h-1)$, divisible by $m$; higher monomials remain in $(z^h)$.
Thus this quotient is $B_{a,e}$-stable and has rank $h$.
For $k=e+2t$ the initial polynomial is $(1+z)^t$, and
\[
 b_{a,k}(n)\equiv[z^0]B_{a,e}^n(1+z)^t\pmod m.
\]
By Cayley--Hamilton, any such scalar sequence satisfies the monic
degree-$h$ characteristic recurrence over $\ZZ/m\ZZ$.
Consequently $h$ consecutive zero differences
$b(n+P)-\gamma b(n)$ imply all subsequent zero differences.

For $(a,e,m)=(8,1,13)$, recurrence \eqref{eq:matrix} gives
\[
 (a(0),\ldots,a(12))\equiv(1,1,3,5,7,8,6,6,5,4,3,9,10)\pmod{13}.
\]
The seven differences $a(j+6)-6a(j)$ vanish for $0\le j\le6$.
Since $h=7$, it follows that
\begin{equation}\label{eq:scalar}
 a(n+6)\equiv6a(n)\pmod{13}\qquad(n\ge0).
\end{equation}
The unit six has order twelve modulo thirteen, so $72$ is a period from
index zero. It is the least period: $a(36)\equiv6^6\equiv12\not\equiv a(0)$ and
$a(24)\equiv6^4\equiv9\not\equiv a(0)\pmod{13}$ exclude shifts $72/2$ and $72/3$.
Every proper divisor of $72$ divides one of these shifts, and the least
period divides every period of a purely periodic sequence.
Equation~\eqref{eq:scalar} also proves
$a(j+6\ell)=6^\ell a(j)\pmod{13}$ for $0\le j<6$ and $\ell\ge0$.
The source's displayed rectangle starts at $a(1)$; its first twelve entries
therefore omit the initial $a(0)=1$.
More generally the characteristic-$p$ identity above gives
$b_{a,k}(n+(p-1)/2)=\alpha_{a,p}b_{a,k}(n)\pmod p$ for $n\ge1$.
It explains the geometric rows at odd primes. The source's precise layouts
also have rank-$h$ scalar certificates: A281180 modulo $7$ has shift $6$
and multiplier $2$ from index $1$; A281181 modulo $13$ has shift $6$ and
multiplier $7$ from index $0$; A281183 and A281184 modulo $11$ have
shift $5$ and multiplier $5$ from index $1$. All five scalar certificates
are included with the supplementary program.

\begin{corollary}[Composite residue rows]\label{cor:rows}
Let $m\ge2$, $r_* =\max_{p\mid m}v_p(m)$, and, for $p^r\Vert m$, set
\[
 h_{p^r}=(p-1)p^{r-1}/2\ (p\text{ odd}),\qquad
 h_2=1,\quad h_{2^r}=2^{r-2}\ (r\ge2).
\]
Put $H=\operatorname{lcm}_{p^r\Vert m}h_{p^r}$.
Choose a unit $C_m$ whose reduction modulo each odd $p$ is
$\alpha_{4,p}^{H/h_{p^r}}$, and whose reduction modulo two is one.
For every $k\ge0$, $n\ge r_*$, and $j\ge0$, the rows
$u_j=b_{4,k}(n+jH)$ satisfy
\[
 \sum_{t=0}^{r_*}(-1)^{r_*-t}\binom{r_*}{t}C_m^{r_*-t}u_{j+t}
 \equiv0\pmod m.
\]
For squarefree $m$, this is the geometric row law $u_{j+1}=C_mu_j$.
\end{corollary}
\begin{proof}
This is Su's Proposition 3.4, with the following direct deduction from the
operators above. Set $A=D^2$ on $M_{p^r}$.
For odd $p$, \eqref{eq:restricted} gives
$(A^{h_{p^r}}-\alpha_{4,p})D=0$ modulo $p$, hence the same identity with
$D$ replaced by $A$. At two use $A^{h_{2^r}+1}=A$ modulo two.
As $h_{p^r}\mid H$, $(A^H-C_m)A$ is divisible by $p$.
Its commuting $r$th power is zero modulo $p^r$.
Multiplication by $(A^H-C_m)^{r_*-r}A^{n-r}$ and evaluation on $c^k$
prove the displayed relation locally; the Chinese remainder theorem combines it.
\end{proof}

For completeness, the same rank-$h$ criterion yields the following
source-specific certificates. The first seven are in Su
\cite[Proposition 3.6]{su}; the last is Bala's A281428 example.
The supplementary program gives every zero difference and every shorter-shift
witness. In the last column, $q:(i;u,v)$ means that shift $P/q$ has
values $u,v$ at index $i$ and $i+P/q$.
\begin{center}\small
\begin{tabular}{lrrrrl}
\toprule
Entry & $m$ & $P$ & Start & $h$ & Shorter-shift witnesses\\
\midrule
A281180 & 7 & 18 & 1 & 4 & $2:(1;1,6),\ 3:(1;1,2)$\\
A281181 & 10 & 8 & 0 & 5 & $2:(0;1,9)$\\
A281181 & 9 & 3 & 1 & 5 & $3:(1;1,4)$\\
A281181 & 13 & 72 & 0 & 7 & $2:(0;1,12),\ 3:(0;1,9)$\\
A281183 & 11 & 25 & 1 & 6 & $5:(1;2,10)$\\
A281184 & 11 & 25 & 1 & 6 & $5:(1;3,4)$\\
A281440 & 13 & 72 & 0 & 7 & $2:(0;1,12),\ 3:(0;1,9)$\\
A281428 & 27 & 18 & 7 & 14 & $2:(7;1,26),\ 3:(7;1,19)$\\
\bottomrule
\end{tabular}
\end{center}
For the mod-$12$ observation attributed to Paul Curtz in A281181, the same
operator with rank six gives $b_{4,1}(j)=1$ for $0\le j\le6$.
Its six successive differences vanish, so $b_{4,1}(n)=1\pmod{12}$ for
every $n\ge0$. This proves the recorded divisibility of every shift difference.
For A281428 modulo $385$, the first six residues are
$(1,1,17,95,315,105)$. The $193$ residues at indices $6$ through $198$
are zero. Since $h=193$, the characteristic recurrence proves that all
terms from index $6$ are zero. At modulus $27$, fourteen zero differences
at indices $7$ through $20$ prove period $18$ from index $7$.
The shorter-shift witnesses in the table prove its least value.
These are the two concrete examples signed by Bala on 3 March 2026,
and are consequences of the same finite-dimensional method.

\section{Mathematical verification scope}
The pinned Lean declaration
\href{https://github.com/the-omega-institute/trureturing/blob/92de86888730265f84e24d28bdb67546c543a966/D5/S1/Recurrence/Periodic/BalaA281440TotientSquarePeriod.lean}{\nolinkurl{BalaA281440TotientSquarePeriod.result}}
proves existence of rational formal series $H,C$, $H(0)=0$, $H'=C^7$,
$C=(\exp H+\exp(-H))/2$, and integer even exponential coefficients with
an eventual shift $\varphi(m)^2$ for every $m\ge1$.
The coefficient-uniqueness argument in Section~1 identifies these witnesses
with the source sequence. In the pinned
\href{https://github.com/the-omega-institute/trureturing/blob/92de86888730265f84e24d28bdb67546c543a966/D5/S1/Recurrence/Periodic/BalaA281181TotientSquarePeriod.lean}{\nolinkurl{BalaA281181TotientSquarePeriod.lean}},
the theorem \texttt{even\_sequence\_period} in the
\texttt{BalaPolynomialSystem} namespace proves the
$\varphi(m)^2$ shift for exponents $a=4,8$ and every power $k\ge0$.
That file's \texttt{result} specializes to A281181.
The three further pinned files
\href{https://github.com/the-omega-institute/trureturing/blob/92de86888730265f84e24d28bdb67546c543a966/D5/S1/Recurrence/Periodic/BalaA281180TotientSquarePeriod.lean}{\nolinkurl{BalaA281180TotientSquarePeriod.lean}},
\href{https://github.com/the-omega-institute/trureturing/blob/92de86888730265f84e24d28bdb67546c543a966/D5/S1/Recurrence/Periodic/BalaA281183TotientSquarePeriod.lean}{\nolinkurl{BalaA281183TotientSquarePeriod.lean}} and
\href{https://github.com/the-omega-institute/trureturing/blob/92de86888730265f84e24d28bdb67546c543a966/D5/S1/Recurrence/Periodic/BalaA281184TotientSquarePeriod.lean}{\nolinkurl{BalaA281184TotientSquarePeriod.lean}}
each contain a \texttt{result} proving the corresponding full-square shift.
The first uses $S'=(1+S^2)^2$ and odd exponential coefficients at indices
$2n-1$, $n\ge1$; the other two use the even exponential coefficients
of $C_4^2$ and $C_4^3$. The last also includes
$(\log(C_4+S_4))'=C_4^3$. All these statements include integrality and
the specified rational-series witnesses.
The sharper half-square period, explicit thresholds, every integer $a\ge2$,
and the derivative and algebraic identities have written proofs above. The least periods
and residue identities have the finite-dimensional proofs of Section~\ref{sec:finite}.
These are distinct scopes.

\paragraph{Supplementary material.}
The accompanying source package contains this \LaTeX{} source, a Python 3
program \texttt{verify.py}, its exact arithmetic certificates in
\texttt{checks.json}, and build instructions. The program compares two
integer coefficient methods at 396 entries across twelve source families,
matches their displayed prefixes (including all 17 A281440 terms), and
supplies 192 uniform-period certificates and 420 further odd-exponent
certificates, in addition to the eight least-period certificates above,
the scalar relations and the algebraic identities.
The universal theorem is established by the written prime-power proof,
independently of these finite checks.

\paragraph{License.}
This note is licensed under Creative Commons Attribution 4.0 International
(CC BY 4.0).

\begin{thebibliography}{9}
\bibitem{su}
Lien-Hung Su, \emph{Quadratic Derivations and Uniform Modular Periods for Six
Inverse-Function Sequences}, version 1, 31 October 2026,
\href{https://doi.org/10.5281/zenodo.23835577}{doi:10.5281/zenodo.23835577}.
Theorem 3.2, Corollary 3.3, Propositions 3.4--3.6 and Theorem 4.1.
\bibitem{oeis440}
OEIS Foundation Inc., \href{https://oeis.org/A281440}{A281440}, revision 17,
6 November 2026. Sequence submitted by Paul D. Hanna, 21 January 2017;
period conjecture by Peter Bala, 19 February 2026.
\bibitem{oeis180}
OEIS Foundation Inc., \href{https://oeis.org/A281180}{A281180}, revision 86,
6 November 2026. Period conjecture by Peter Bala, 18 February 2026.
\bibitem{oeis181}
OEIS Foundation Inc., \href{https://oeis.org/A281181}{A281181}, revision 103,
6 November 2026. Period conjecture by Peter Bala, 16 February 2026;
residue observations by Paul Curtz, 20 January 2017;
integer column recurrence by Seiichi Manyama, 14 April 2026.
\bibitem{oeis183}
OEIS Foundation Inc., \href{https://oeis.org/A281183}{A281183}, revision 60,
6 November 2026. Period conjecture by Peter Bala, 16 February 2026.
\bibitem{oeis184}
OEIS Foundation Inc., \href{https://oeis.org/A281184}{A281184}, revision 41,
6 November 2026. Period conjecture by Peter Bala, 16 February 2026.
\bibitem{oeiscompanions}
OEIS Foundation Inc., companion entries
\href{https://oeis.org/A281427}{A281427} (rev. 29),
\href{https://oeis.org/A281428}{A281428} (rev. 54),
\href{https://oeis.org/A281432}{A281432} (rev. 10),
\href{https://oeis.org/A281436}{A281436} (rev. 10),
\href{https://oeis.org/A281439}{A281439} (rev. 6),
\href{https://oeis.org/A281443}{A281443} (rev. 5),
\href{https://oeis.org/A281444}{A281444} (rev. 9).
All OEIS entries consulted 14 November 2026.
\bibitem{dominici}
Diego Dominici, \emph{Nested derivatives: a simple method for computing
series expansions of inverse functions}, International Journal of Mathematics
and Mathematical Sciences \textbf{2003} (58), 3699--3715,
\href{https://doi.org/10.1155/S0161171203303291}{doi:10.1155/S0161171203303291};
\href{https://arxiv.org/abs/math/0501052v2}{arXiv:math/0501052v2}, Theorem 4.1.
\bibitem{graves}
Peter Bala, \emph{The method of Graves for computing inverse functions},
26 January 2026,
\url{https://oeis.org/A281181/a281181.pdf}, Corollary 1.
\end{thebibliography}
\end{document}
